\section{Arquitetura do modelo}
\label{sec:arquitetura}

Foi utilizada a EfficientNet-B0~\cite{tan2019efficientnet}, na implementação
da biblioteca Torchvision, com pesos pré-treinados no ImageNet. O extrator
de características contém 16 blocos MBConv, que combinam convoluções
\textit{depthwise} e pontuais, Batch Normalization e mecanismos
\textit{squeeze-and-excitation} para recalibrar os canais. A rede utiliza
ativações SiLU e conexões residuais nos blocos com dimensões compatíveis.

Após a extração de características, uma operação de média global espacial
produz um vetor com 1.280 componentes por imagem. O classificador contém
Dropout com probabilidade $0{,}2$ e uma camada linear. A camada original,
com 1.000 saídas, foi substituída por uma nova camada com 101 saídas,
inicializada aleatoriamente. Essa camada possui
$1.280 \times 101 + 101 = 129.381$ parâmetros, incluindo os vieses.

O modelo adaptado possui \textbf{4.136.929 parâmetros}, dos quais 4.007.548
pertencem ao extrator e 129.381 ao classificador. A
Tabela~\ref{tab:parametros} apresenta a quantidade treinável em cada
experimento. A liberação dos parâmetros de escala e deslocamento do Batch
Normalization acrescenta 42.016 parâmetros ao conjunto otimizado. Suas
médias e variâncias acumuladas são estatísticas armazenadas pelo modelo e
não entram nessa contagem. O Stochastic Depth também não adiciona parâmetros.

\begin{table}[ht]
  \centering
  \caption{Quantidade de parâmetros do modelo adaptado por experimento.}
  \label{tab:parametros}
  \begin{tabular}{lrr}
    \hline
    Experimento & Treináveis & Congelados \\
    \hline
    Modelo de referência & 129.381 & 4.007.548 \\
    BN estatístico & 129.381 & 4.007.548 \\
    BN treinável & 171.397 & 3.965.532 \\
    BN treinável + SD & 171.397 & 3.965.532 \\
    \hline
  \end{tabular}
\end{table}
